\section{What is optimized and what is proved}
The object sought is a fixed polynomial with integer coefficients, not a
family with an uninstantiated machine parameter. Its quantified variables
have one explicitly stated domain. A polynomial abbreviation and a shared
arithmetic register introduce no extra variables; conversely, a quantified
witness hidden inside a special-function graph must be counted.

We compare the number of quantified variables, exact total degree, and
live binary arithmetic operations. Fixed integer constants are coefficients,
not additional unknowns. Addition, subtraction and multiplication each cost
one operation, including multiplication by a small constant. This convention
measures a shared straight-line circuit, not expanded monomial count or bit
complexity. Large fixed coefficients remain an important separate cost.
There is no assertion that these measures admit a common global minimum.

The new results are ordinary mathematical reductions. Their concluding
consistency statements compose these reductions with explicitly identified,
pinned arithmetic-kernel, initialized-machine and syntactic-refutation
calibration theorems. The accompanying programs check literal circuits,
identities, degree certificates and fixed coefficient payloads. They do not
replace the unbounded mathematical proofs, and do not constitute a new
end-to-end Lean, Rocq or encoded object-language ZFC certificate.

In particular, the assertion is an equivalence
\[
 \mathsf{ZFC}\vdash
 \bigl(\neg\exists\mathbf x\in D^n\ P(\mathbf x)=0\bigr)
 \longleftrightarrow\operatorname{Con}(\mathsf{ZFC}),
\]
using the stated ordinary calibration. It is not a proof of either side.
Here $D=\mathbb Z$, $\mathbb Z_{>0}$ or $\mathbb N$ is specified separately
for each construction, and $\mathbb N$ includes zero.

\section{Coprime products and monic congruence carriers}
The elementary observation driving this iteration is stronger when used
in the design of witnesses rather than applied after a polynomial has
already been built.

\begin{lemma}[Coprime square extraction]
Let $a_1,\ldots,a_m$ be positive, pairwise coprime integers. Their product
is a square if and only if each $a_i$ is a square. If $a,b,d$ are positive integers, $ab/d$ is a rational square and
$\gcd(a,bd)=1$, then $a$ is a square.
\end{lemma}
\begin{proof}
For every prime dividing $a_i$, the valuation of the product is exactly
the valuation of $a_i$. It must be even. The rational assertion uses the
same valuation computation, with the denominator contributing zero at
these primes. Positivity fixes the sign ambiguity.
\end{proof}

For example, $a$ and $1+ad u^2$ are automatically coprime. Imposing that
their product is a square therefore recovers two square conditions from
one. This is not a universal equivalence between arbitrary two square
conditions and one product condition: the coprimality and sign proofs are
essential. It also changes the original witness problem. Completeness must
show that the newly constrained Pell ordinate can actually be chosen.

\begin{lemma}[A useful two-root criterion]
Let $a,b>0$ be rational numbers with $a\ne b$. If a signed sum of their
square roots is rational, both square roots are rational.
\end{lemma}
\begin{proof}
The sum cannot be zero. If $z=\epsilon\sqrt a+\eta\sqrt b\ne0$, then
$2z\epsilon\sqrt a=z^2+a-b$, and similarly for the other root.
\end{proof}

Neither statement can be used without its hypotheses. The integer
constructions below prove positivity before splitting products. The
positive-variable construction proves the two radicands unequal by a
congruence obstruction, before making any parity or transport deduction.

\begin{lemma}[A monic carrier for two congruences]
Let $M,H$ be nonzero coprime integers, and let $n,K,v,v_0$ be integers.
Set $U=MH$ and $Z=Un+MK+H(v-v_0)$. If $Z/U$ is a root of any monic
polynomial with integer coefficients, then
\[
 M\mid v-v_0,\qquad H\mid K.
\]
\end{lemma}
\begin{proof}
A rational algebraic integer is an integer, so $U\mid Z$. Reducing modulo
$M$ and $H$ and using coprimality gives the two conclusions.
\end{proof}

This observation changes the order of the proof. A monic radical norm
can force the scale and projection congruences before any individual
radical has been shown rational. Its norm equation alone makes the total
rational center integral. The restored scale can then justify the bounds
or prime-valuation arguments needed to recover each separate Pell root.
This avoids a circular appeal to the very scale that the equation is
supposed to enforce.
